\chapter{ARDC in practice}
\label{ch:ardc}

This is the operational core of the note: \textbf{A}rchive, \textbf{R}eference,
\textbf{D}escribe, \textbf{C}ite \& credit. Each step has concrete tools, and the chapter is
punctuated by hands-on labs you run on a repository of your choice --- your own code, or a
report, dependency, or library you want to mention, cite, or rely upon.

\takeaway{Archive once, then reference at the granularity your claim needs.}

\section{Archive}
% TODO draft from: swh-acquisition-process.org::#savecodenow; deposit.org::#overview,#prepare,#send,#status.
Save Code Now, the browser extension, a release webhook, or the SWORD deposit API.

\begin{handson}[Hands-on: archive any repository and get its SWHID]
\begin{enumerate}[nosep,leftmargin=*]
  \item Pick a repository worth referencing: your own, or one you depend on / want to cite.
  \item Archive it --- either with the \textbf{updateswh} browser extension (one click on a
        GitHub/GitLab/Bitbucket page) or via \url{https://save.softwareheritage.org}.
  \item Wait for the visit to reach status \texttt{full}; open the object and copy the
        \emph{qualified} \textsc{swhid} from the \emph{Permalinks} tab.
  \item Confirm it resolves under \texttt{https://archive.softwareheritage.org/}.
\end{enumerate}
\end{handson}

% Hands-on (to draft): set up a release WEBHOOK so each tag is auto-archived ---
%   GitHub, self-hosted GitLab, Gitea, Bitbucket, SourceForge.

\section{Reference}
% TODO draft from: swh-ardc.org::#swh-r. Granularity snp/rev/rel/dir/cnt/lines;
% qualifiers origin=/visit=/anchor=; worked examples (appendix B).

% Hands-on (to draft): build a reference to embed in a research article ---
%   a code snippet (cnt + lines), a revision (rev), a file, and a directory ---
%   and paste each into the archive to confirm it highlights the right object.

\section{Describe}
% TODO draft from: research-software-essentials.org::#description,#licence,#versioncontrol.
% codemeta.json, CITATION.cff, SPDX, REUSE.

% Hands-on (to draft): create a codemeta.json, commit it into your repo, re-archive,
%   then check how the citation sidebar appears in the Archive and use it to produce
%   the citation(s) you want.

\section{Cite and credit}
% TODO draft from: swh-ardc.org::#swh-dc. Refs: DiCosmoSEN2020, cise-2020-doi,
% force11citationprinciples, SCIDWG2020.
The \texttt{biblatex-software} style adds four entry types and a \texttt{swhid=}
field~\autocite{DiCosmoSEN2020,cise-2020-doi}. For example, this note cites a specific
version of a library, \autocite{parmap-1.1.1}, and a code fragment within
it, \autocite{simplemapper}.

\begin{handson}[Hands-on: cite software with biblatex-software]
\begin{enumerate}[nosep,leftmargin=*]
  \item Write a \texttt{@softwareversion} entry for the code you archived, with its
        \texttt{swhid=} field; then add a \texttt{@codefragment} child via \texttt{crossref}.
  \item Produce \emph{variants}: a project-level \texttt{@software}, a \texttt{@softwaremodule},
        and a line-pinpointed fragment --- and compare how each renders.
  \item Enable \texttt{biblatex-software} inside an \texttt{acmart.cls} article and confirm the
        \textsc{swhid} renders in the ACM template.
\end{enumerate}
\end{handson}
